% !TEX program = xelatex
% =====================================================================
%  colorbeam lecture template
%  Compile with XeLaTeX or LuaLaTeX, twice:  xelatex colorbeam-template.tex
%  Required files in the same folder: beamerthemecolorbeam.sty and the
%  colorbeamlib-*.sty libraries loaded below
%  Home: https://github.com/eskander/colorbeam
% =====================================================================
% ---------------------------------------------------------------------
%  Copyright 2026 Eskander Bejaoui
%
%  This work may be distributed and/or modified under the
%  conditions of the LaTeX Project Public License, either version 1.3c
%  of this license or (at your option) any later version.
%  The latest version of this license is in
%    https://www.latex-project.org/lppl.txt
%  and version 1.3c or later is part of all distributions of LaTeX
%  version 2008 or later.
%
%  This work has the LPPL maintenance status `maintained'.
%
%  The Current Maintainer of this work is Eskander Bejaoui.
%  Bug reports and contributions: https://github.com/eskander/colorbeam
%
%  This work consists of the files beamerthemecolorbeam.sty,
%  colorbeamlib-programming.sty, colorbeamlib-electronics.sty,
%  colorbeam-template.tex, README.md and LICENSE.
% ---------------------------------------------------------------------
\documentclass[aspectratio=169,12pt,t]{beamer}

% Fonts (optional, before \usetheme). Default: Carlito + DejaVu Sans Mono.
% \def\colorbeamFont{Source Sans 3}
% \def\colorbeamMono{Fira Code}
\usetheme{colorbeam}

% Optional subject libraries, load only what you need (each one pulls in extra packages):
%   programming : Arduino, Python, C, C++ and JavaScript code blocks
%   electronics : CircuiTikZ, \ohm and board/wire diagram styles
\usecolorbeamlibrary{programming,electronics}

% Main color of the deck (title slide, introduction, \resetpart).
% Any palette color works: grape, sky, mint, leaf, sun, clay, coral, rose, plum, slate.
% \setmaincolor{plum}

% Language: polyglossia (XeLaTeX). For another language, change the main language
% and the word used in "Part 1", e.g. \setmainlanguage{french} and \renewcommand{\partlabel}{Partie}
\usepackage{polyglossia}
\setmainlanguage{english}

% Small helper used on the library slides of this template
\newcommand{\libnote}[1]{\pill{slate}{library: #1}\hspace{0.6em}{\scriptsize\color{ink!70}\texttt{\textbackslash usecolorbeamlibrary\{#1\}}}\par\smallskip}

% Footer text
\lecturefooter{My course \textbullet\ Session 1}

\begin{document}

% ---------------------------------------------------------------------
% Title: \titleslide[color]{kicker}{title}{subtitle}{subtitle 2}{bottom note}
% ---------------------------------------------------------------------
\titleslide{COURSE \textbullet\ MY SUBJECT}{Session title}{A clear subtitle}{possibly on two lines}{First Last \textbullet\ University \textbullet\ 2026}

% ---------------------------------------------------------------------
% Outline with numbered steps (each step can have its own color)
% ---------------------------------------------------------------------
\begin{frame}{Session outline}
  \begin{columns}[T,onlytextwidth]
    \begin{column}{0.6\textwidth}
      \step{mint}{1}{7.3cm}Boxes and callouts\estep
      \step{sky}{2}{7.3cm}Tables and code, with a longer line to show the automatic line wrapping\estep
      \step{leaf}{3}{7.3cm}Diagrams, circuits and plots\estep
      \step{plum}{4}{7.3cm}Questions and wrap-up\estep
    \end{column}
    \begin{column}{0.4\textwidth}
      \begin{card}
        {\bfseries\color{accent}Goal}\par\smallskip
        One sentence that sums up what students will be able to do at the end.
      \end{card}
    \end{column}
  \end{columns}
\end{frame}

% ---------------------------------------------------------------------
% The palette: ten colors, each usable as a section color
% ---------------------------------------------------------------------
\begin{frame}{The colorbeam palette}
  \centering
  \begin{tikzpicture}
    \foreach \c/\hex/\nm/\role [count=\i from 0] in {
        grape/2A3F73/midnight blue/{default, intro},
        sky/2F6DA3/steel blue/{section, information},
        mint/1D7874/teal/{section},
        leaf/4F8A3C/green/{section, success},
        sun/C28A12/amber/{tips, highlights},
        clay/B8622B/terracotta/{section},
        coral/B5443B/muted red/{warnings, questions},
        rose/A8416B/raspberry/{section},
        plum/6A4A9C/violet/{section},
        slate/55657A/slate gray/{neutral, secondary}}
      {
        \pgfmathsetmacro{\px}{mod(\i,5)*2.85}
        \pgfmathsetmacro{\py}{-floor(\i/5)*3.05}
        \fill[\c,rounded corners=6pt] (\px,\py) rectangle ++(2.55,-1.45);
        \node[anchor=north west,text=white,font=\large\bfseries] at (\px+0.15,\py-0.1) {\c};
        \node[anchor=south west,text=white!85,font=\scriptsize] at (\px+0.15,\py-1.35) {\#\hex};
        \node[anchor=north,text=ink,font=\footnotesize\bfseries,text width=2.5cm,align=center] at (\px+1.275,\py-1.52) {\nm};
        \node[anchor=north,text=ink!65,font=\scriptsize,text width=2.5cm,align=center] at (\px+1.275,\py-1.88) {\role};
      }
  \end{tikzpicture}
\end{frame}

% ---------------------------------------------------------------------
% Pills and colored steps
% ---------------------------------------------------------------------
\begin{frame}{Pills and numbered steps}
  \begin{columns}[T,onlytextwidth]
    \begin{column}{0.5\textwidth}
      {\bfseries\color{accent}Steps in any color}\par\medskip
      \step{grape}{1}{5.5cm}Plan\estep
      \step{sky}{2}{5.5cm}Build\estep
      \step{leaf}{3}{5.5cm}Test\estep
      \step{sun}{4}{5.5cm}Fix\estep
      \step{clay}{5}{5.5cm}Document\estep
    \end{column}
    \begin{column}{0.5\textwidth}
      {\bfseries\color{accent}Pills}\par\medskip
      \pill{grape}{grape} \pill{sky}{sky} \pill{mint}{mint} \pill{leaf}{leaf}\par\medskip
      \pill{sun}{sun} \pill{clay}{clay} \pill{coral}{coral} \pill{rose}{rose}\par\medskip
      \pill{plum}{plum} \pill{slate}{slate}\par\bigskip
      \begin{card}\small\raggedright
        An empty color, as in \texttt{\textbackslash step\{\}\{1\}\{5cm\}}, uses the color of the current section.
      \end{card}
    \end{column}
  \end{columns}
\end{frame}

% ---------------------------------------------------------------------
% Part slide: changes the color and label of the following slides
% \partslide{number}{color}{title}{subtitle}
% ---------------------------------------------------------------------
\partslide{1}{mint}{Boxes and callouts}{highlight what matters}

\begin{frame}{Callout boxes}
  \small
  \tcbset{top=1.4mm,bottom=1.4mm}
  \begin{columns}[T,onlytextwidth]
    \begin{column}{0.5\textwidth}
      \begin{card}
        {\bfseries\color{accent}card}\par\smallskip
        Neutral box, section color.
      \end{card}
      \vspace{0.2cm}
      \begin{tip}
        {\bfseries tip}\par\smallskip
        A practical remark.
      \end{tip}
      \vspace{0.2cm}
      \begin{infobox}
        {\bfseries infobox}\par\smallskip
        Extra information or a reminder.
      \end{infobox}
    \end{column}
    \begin{column}{0.5\textwidth}
      \begin{warnbox}
        {\bfseries warnbox}\par\smallskip
        A warning: do not miss this!
      \end{warnbox}
      \vspace{0.2cm}
      \begin{okbox}
        {\bfseries okbox}\par\smallskip
        A result or a good practice.
      \end{okbox}
      \vspace{0.2cm}
      \begin{itemize}
        \item Plain bullet list
        \item Bullets follow the section color
      \end{itemize}
    \end{column}
  \end{columns}
\end{frame}

\begin{frame}{Any color, any box}
  \begin{columns}[T,onlytextwidth]
    \begin{column}{0.33\textwidth}
      \begin{colorcard}{clay}
        {\bfseries\color{clay}colorcard\{clay\}}\par\smallskip\small
        Tinted background, colored frame.
      \end{colorcard}
    \end{column}
    \begin{column}{0.33\textwidth}
      \begin{colorcard}{rose}
        {\bfseries\color{rose}colorcard\{rose\}}\par\smallskip\small
        Use it for a recurring theme.
      \end{colorcard}
    \end{column}
    \begin{column}{0.33\textwidth}
      \begin{colorcard}{plum}
        {\bfseries\color{plum}colorcard\{plum\}}\par\smallskip\small
        Or to tell two ideas apart.
      \end{colorcard}
    \end{column}
  \end{columns}
  \vspace{0.3cm}
  \begin{columns}[T,onlytextwidth]
    \begin{column}{0.5\textwidth}
      \begin{colorcard}{slate}
        {\bfseries\color{slate}colorcard\{slate\}}\par\smallskip\small
        Neutral side notes and references.
      \end{colorcard}
    \end{column}
    \begin{column}{0.5\textwidth}
      \begin{colorcard}{leaf}
        {\bfseries\color{leaf}colorcard\{leaf\}}\par\smallskip\small
        Same look as okbox, any color you like.
      \end{colorcard}
    \end{column}
  \end{columns}
  \vspace{0.35cm}
  \begin{darkbox}
    {\bfseries darkbox: the one thing to remember}
  \end{darkbox}
\end{frame}

% ---------------------------------------------------------------------
\partslide{2}{sky}{Tables and code}{core, plus the programming library}

\begin{frame}{Comparison table}
  \renewcommand{\arraystretch}{1.3}
  \centering
  \begin{tabular}{@{}>{\bfseries}p{4cm}p{4.2cm}p{4.2cm}@{}}
    \rowcolor{ink}\color{white}\bfseries & \color{white}\bfseries Option A & \color{white}\bfseries Option B\\
    \rowcolor{white} Criterion 1 & value & value\\
    \rowcolor{accent!10} Criterion 2 & value & \textcolor{leaf}{\textbf{best value}}\\
    \rowcolor{white} Criterion 3 & value & value\\
    \rowcolor{accent!10} Criterion 4 & \textcolor{coral}{\textbf{worst value}} & value\\
  \end{tabular}
\end{frame}

% Frames that contain code must be [fragile]
\begin{frame}[fragile]{Code block (core)}
  \begin{columns}[T,onlytextwidth]
    \begin{column}{0.62\textwidth}
      \begin{codebox}[title={example.py},listing options={style=code,language=Python}]
def average(values):
    # compute the mean
    return sum(values) / len(values)

print(average([12, 15, 9]))
      \end{codebox}
    \end{column}
    \begin{column}{0.38\textwidth}
      \begin{card}
        {\bfseries\color{accent}Your turn}\par\smallskip\small
        \step{}{1}{3.2cm}Write the code\estep
        \step{}{2}{3.2cm}Run it\estep
        \step{}{3}{3.2cm}Compare the result\estep
      \end{card}
    \end{column}
  \end{columns}
\end{frame}

\begin{frame}[fragile]{Python and JavaScript}
  \libnote{programming}
  \begin{columns}[T,onlytextwidth]
    \begin{column}{0.48\textwidth}
      \begin{pycode}[title={mean.py},listing options={style=py,basicstyle=\codefont{9}{11}}]
def mean(values):
    # arithmetic mean
    return sum(values) / len(values)

data = [12, 15, 9, 14]
print(mean(data))
      \end{pycode}
    \end{column}
    \begin{column}{0.52\textwidth}
      \begin{jscode}[title={mean.js},listing options={style=js,basicstyle=\codefont{9}{11}}]
const data = [12, 15, 9, 14];

function mean(v) {
  // arithmetic mean
  const s = v.reduce((a, b) => a+b);
  return s / v.length;
}

console.log(mean(data));
      \end{jscode}
    \end{column}
  \end{columns}
\end{frame}

\begin{frame}[fragile]{C and C++}
  \libnote{programming}
  \begin{columns}[T,onlytextwidth]
    \begin{column}{0.48\textwidth}
      \begin{ccode}[title={mean.c},listing options={style=c,basicstyle=\codefont{9}{11}}]
#include <stdio.h>

int main(void) {
  int v[] = {12, 15, 9, 14};
  int s = 0;
  for (int i = 0; i < 4; i++)
    s += v[i];  /* sum */
  printf("%.1f\n", s / 4.0);
  return 0;
}
      \end{ccode}
    \end{column}
    \begin{column}{0.52\textwidth}
      \begin{cppcode}[title={mean.cpp},listing options={style=cpp,basicstyle=\codefont{9}{11}}]
#include <iostream>
#include <vector>

int main() {
  std::vector<int> v{12, 15, 9, 14};
  double s = 0;
  for (int x : v) s += x;
  std::cout << s / v.size()
            << std::endl;
}
      \end{cppcode}
    \end{column}
  \end{columns}
\end{frame}

\begin{frame}[fragile]{Arduino, two blocks side by side}
  \libnote{programming}
  \begin{columns}[T,onlytextwidth]
    \begin{column}{0.42\textwidth}
      \begin{ino}[title={setup},listing options={style=ino,basicstyle=\codefont{9.5}{11.6}}]
const int LED = 2;

void setup() {
  pinMode(LED, OUTPUT);
  Serial.begin(115200);
}
      \end{ino}
    \end{column}
    \begin{column}{0.58\textwidth}
      \begin{ino}[title={loop},listing options={style=ino,basicstyle=\codefont{9.5}{11.6}}]
void loop() {
  digitalWrite(LED, HIGH);  // on
  delay(500);
  digitalWrite(LED, LOW);   // off
  delay(500);
}
      \end{ino}
    \end{column}
  \end{columns}
\end{frame}

% ---------------------------------------------------------------------
\partslide{3}{leaf}{Diagrams and plots}{TikZ, pgfplots and the electronics library}

\begin{frame}{Chained formulas with flowbox}
  \centering
  \begin{tikzpicture}
    \node[flowbox]      (a) at (0,0)   {{\bfseries\color{accent}1. Measure}\\[6pt] $V=\dfrac{N\times V_{\mathrm{ref}}}{N_{\max}}$};
    \node[flowbox=sky]  (b) at (4.6,0) {{\bfseries\color{sky}2. Convert}\\[6pt] $R=\dfrac{R_0\,V}{V_{\mathrm{ref}}-V}$};
    \node[flowbox=plum] (c) at (9.2,0) {{\bfseries\color{plum}3. Result}\\[6pt] $T=f(R)$};
    \draw[ar] (a.east) -- (b.west);
    \draw[ar] (b.east) -- (c.west);
  \end{tikzpicture}
  \vspace{0.4cm}

  {\small\color{ink!70}\texttt{flowbox} takes an optional color: \texttt{flowbox=plum}. Without it, it uses the section color.}
\end{frame}

\begin{frame}{Electrical schematic}
  \libnote{electronics}
  \begin{columns}[T,onlytextwidth]
    \begin{column}{0.6\textwidth}
      \begin{circuitikz}[line width=1.3pt,every node/.style={font=\footnotesize}]
        \draw (1.9,4.3) to[R,l_=$10\,\mathrm{k}\ohm$] (1.9,2.4);
        \draw (1.9,2.4) to[R,l_=$10\,\mathrm{k}\ohm$] (1.9,0.5);
        \draw[wire,coral] (1.9,4.3) -- (5.4,4.3);
        \draw[wire,black!80] (1.9,0.5) -- (5.4,0.5);
        \draw[wire,orange!90!black] (1.9,2.4) node[circ]{} -- (5.4,2.4);
        \node[board,minimum width=2.7cm,minimum height=4.9cm,anchor=south west] at (5.4,0.05) {};
        \foreach \y/\l in {4.3/3V3,2.4/GPIO 34,0.5/GND} {\fill[pinc] (5.4,\y) circle (0.1); \node[pinl] at (5.6,\y) {\l};}
        \node[rotate=90,text=white,font=\small\bfseries] at (7.75,2.5) {BOARD};
      \end{circuitikz}
    \end{column}
    \begin{column}{0.4\textwidth}
      \begin{card}
        Voltage divider: the middle point is wired to the analog input.
      \end{card}
      \vspace{0.2cm}
      \begin{tip}\small Red: power. Black: ground. Orange: signal.\end{tip}
    \end{column}
  \end{columns}
\end{frame}

\begin{frame}{RC charging circuit}
  \libnote{electronics}
  \begin{columns}[T,onlytextwidth]
    \begin{column}{0.5\textwidth}
      \begin{circuitikz}[line width=1.3pt,every node/.style={font=\footnotesize}]
        \draw (0,0) to[V,l_=$V_s$] (0,3)
              to[R,l=$R$] (3.6,3)
              to[C,l=$C$] (3.6,0) -- (0,0);
      \end{circuitikz}
    \end{column}
    \begin{column}{0.5\textwidth}
      \begin{card}
        {\bfseries\color{accent}Time constant}\par\smallskip
        $\tau = RC$, for example $10\,\mathrm{k}\ohm \times 100\,\mu\mathrm{F} = 1\,\mathrm{s}$.
      \end{card}
      \vspace{0.2cm}
      \begin{tip}\small The next slide plots three values of $\tau$.\end{tip}
    \end{column}
  \end{columns}
\end{frame}

\begin{frame}{Plot with several colors}
  \centering
  \begin{tikzpicture}
    \begin{axis}[width=11cm,height=5.8cm,xlabel={Time (s)},ylabel={Voltage (V)},
        label style={font=\footnotesize,text=ink},tick label style={font=\scriptsize},axis line style={ink},
        grid=both,grid style={ink!10},every axis plot/.append style={line width=2.2pt},
        legend style={font=\scriptsize,at={(0.97,0.05)},anchor=south east,draw=ink!30}]
      \addplot[sky,domain=0:5,samples=60]{3.3*(1-exp(-x))};
      \addplot[coral,domain=0:5,samples=60]{3.3*(1-exp(-2*x))};
      \addplot[plum,domain=0:5,samples=60]{3.3*(1-exp(-0.5*x))};
      \legend{$\tau=1$ s,$\tau=0.5$ s,$\tau=2$ s}
    \end{axis}
  \end{tikzpicture}
\end{frame}

% ---------------------------------------------------------------------
\partslide{4}{plum}{Questions and wrap-up}{engage the class}

\begin{frame}{Question}
  \questionmark
  \vspace{0.3cm}
  \begin{card}[width=10.5cm,colframe=coral,boxrule=1.6pt,top=5mm,bottom=4mm]
    {\fontsize{22}{28}\selectfont\bfseries A question to ask the class?}\par\bigskip
    {\fontsize{22}{28}\selectfont\bfseries\color{coral}Why?}
  \end{card}
  \vspace{0.5cm}
  {\small\color{ink!70}Discuss with your neighbor before moving on.}
\end{frame}

\begin{frame}{Numbered rows and troubleshooting rows}
  \numrow{}{1}{Numbered row}{A title and a one-line description, framed in the section color}
  \numrow{sun!90!black}{2}{Numbered row in another color}{Pass any palette color as first argument}
  \vspace{0.3cm}
  \helprow{}{Problem}{Short explanation of the cause and of the fix}
  \helprow{coral}{Another problem}{Same row, in red for something critical}
\end{frame}

% ---------------------------------------------------------------------
% Wrap-up: \resetpart returns to the main color, with no label
% ---------------------------------------------------------------------
\resetpart
\begin{frame}{Summary}
  \vspace{0.2cm}
  \step{}{1}{12.2cm}First takeaway\estep
  \step{}{2}{12.2cm}Second takeaway\estep
  \step{}{3}{12.2cm}Third takeaway\estep
  \vspace{0.1cm}
  \begin{darkbox}
    {\bfseries Going further: an idea for the next session}
  \end{darkbox}
\end{frame}

\end{document}
